\section{Introduction}
\label{sec:intro}

Suppose you run a reporting query on \pg and it returns a row that should not
be there. The query does not crash and no error is raised. Nothing in the
logs says that the answer is wrong. This kind of failure, a \emph{wrong
result}, is the most harmful bug a database system can have, because the user
has no way to notice it.

Fuzzers for database systems mostly look for crashes. SQLsmith~\cite{sqlsmith}
generates random queries from the system catalog, and
SQUIRREL~\cite{zhong2020squirrel} mutates SQL with coverage feedback. Both
detect a bug when the server crashes or an assertion fails. To find wrong
results, a fuzzer needs a \emph{test oracle} that knows the right answer.
Differential testing compares two systems~\cite{slutz1998rags}; metamorphic
oracles such as Ternary Logic Partitioning (TLP)~\cite{rigger2020tlp},
NoREC~\cite{rigger2020norec} and equivalent expression
transformation~\cite{jiang2024eet} compare a query with a rewritten query that
must return the same rows. These oracles rewrite the SQL text of a query,
so each tool builds its own small query generator that the oracle understands.

We focus on a gap between the two lines of work. Grammar-based generators
reach deep into the planner but do not know which parts of a query can be
rewritten safely. Oracle-based tools know how to rewrite, but their
generators cover a small part of SQL. And neither knows what changed in the
code yesterday: the planner of \pg master changes every week, and most new
bugs live in the newest code.

Figure~\ref{fig:motivating} shows a real wrong result that \tool found. Since
\pg~17, the planner replaces \texttt{range @> elem} by the constant
\texttt{true} when the range is unbounded on both sides, and by
\texttt{false} when it is empty. The operator is strict, so for a NULL
element the correct answer is NULL. The constant is wrong under \texttt{NOT},
in \texttt{IS NULL} tests, in \texttt{CASE}, and in outer joins, where unmatched
rows carry NULLs. We reproduced the bug on the release 18.6 and on master.
\tool found it with a typed rewrite of the derivation tree (a scalar
subquery around the range), in the first 150 test cases the rewrite ran.

This is not a single case. Between 31~July and 16~October~2026, \pg master
received 23 fixes in query processing that came with a regression test and
still revert cleanly (Section~\ref{sec:rq3}). When we put each bug back, 20 of
them change what a query returns: 8 return wrong rows, 5 raise internal
errors, 4 fail an assertion, and 3 raise or drop an error message. Only 9 of
these 20 can be seen by a crash-only fuzzer.

\begin{figure}[t]
\begin{lstlisting}[numbers=none,escapechar=!,language=SQL]
SELECT c,
       '(,)'::int4range @> c       AS contains,   -- must be NULL when c is NULL
       NOT ('empty'::int4range @> c) AS not_empty  -- must be NULL when c is NULL
FROM (VALUES (NULL::int), (1)) v(c);

  c   | contains | not_empty
------+----------+----------
 NULL | t        | t          !\textbf{$\leftarrow$ wrong: the planner folded the call to a constant}!
 1    | t        | t
\end{lstlisting}
\caption{A wrong result in \pg~17, 18.6 and master that \tool found: the
planner support function for range containment returns \texttt{true} or
\texttt{false} for a NULL element, where the strict operator gives NULL. No
error or crash shows it; only an oracle that compares two equivalent queries
does.}
\label{fig:motivating}
\end{figure}

Closest to us is query plan guidance~\cite{ba2023qpg}, which steers a
generator toward new plan shapes, and Fuzz4All~\cite{xia2024fuzz4all}, which
uses a large language model (LLM) as the generator. However, plan guidance
uses the oracles of SQLancer and their small query language, and an LLM that
writes every test case is slow and produces many invalid queries. Neither
connects the generator to the code that changed.

We present \tool, a grammar-aware fuzzer for \pg built on three ideas:
\begin{enumerate}
  \item A \emph{typed grammar} generates SQL as derivation trees. Every node
    knows its type and its place in the query (Section~\ref{sec:grammar}).
  \item \emph{Tree-based oracles} rewrite the derivation tree, not the text:
    TLP on \texttt{WHERE} and \texttt{HAVING} at any query level, NoREC, ten
    typed expression rewrites, CTE materialisation and table fences. A
    plan-variant oracle and a cross-version oracle compare the same query
    under other planner settings and on the last release
    (Section~\ref{sec:oracles}).
  \item \emph{Commit-directed grammar extension}: an LLM reads a recent commit
    and proposes grammar alternatives and seed queries; a \gate accepts them
    only if the server runs them and they execute the lines the commit
    changed (Section~\ref{sec:llm}).
\end{enumerate}
\tool then reduces every finding to a short reproducer and records a verdict
for it (Section~\ref{sec:triage}).

\num{First results.} In a four-minute smoke run on a server with the 23 bugs
put back, \tool rediscovered an assertion failure in \texttt{JSON\_TABLE}
that we had found earlier on master, and its triage exposed two false-positive
classes in our own oracles, which we fixed (Section~\ref{sec:rq2}).
\num{RQ1--RQ4 numbers follow from the full campaigns.} The lesson so far is
that oracle precision, not generator power, limits a wrong-result fuzzer: a
generator that reaches more of SQL also reaches more corner cases where two
``equivalent'' queries are not equivalent.

\paragraph{Contributions.}
\begin{itemize}
  \item A typed SQL grammar for \pg with node spans, so that oracles can
    rewrite derivation trees (Section~\ref{sec:approach}).
  \item Six metamorphic oracles on derivation trees, plus plan-variant and
    cross-version oracles, with the precision of each (Section~\ref{sec:rq2}).
  \item A commit-directed LLM extension with a \gate that checks validity,
    reachability and diff coverage on the server (Section~\ref{sec:llm}).
  \item A benchmark of 23 real \pg bugs put back into one server, with a
    single-bug server per bug for attribution (Section~\ref{sec:rq3}).
  \item Two previously unknown \pg bugs, one wrong result and one assertion
    failure, with patches (\num{upstream status}; Section~\ref{sec:rq2}), and
    the tool and all data (Section~\ref{sec:data}).
\end{itemize}
